Une solution d'acide chlorhydrique a une concentration molaire
$c=\qty{2.50e-2}{\mol\per\L}$.
\begin{enumerate}
    \item Écrire l'équation de la réaction de dissolution du chlorure
          d'hydrogène dans l'eau.
    \item Calculer les concentrations $\left[\ce{H3O+}\right]$ et
          $\left[\ce{Cl-(aq)}\right]$ dans la solution.
    \item Comment met-on en évidence les ions présents en solution~?
\end{enumerate}

